% Luc Destombe --- licence CC BY-NC-SA 4.0
% https://creativecommons.org/licenses/by-nc-sa/4.0/deed.fr
% Fichier autonome : compiler avec pdflatex, deux fois (signets, nombre de pages).
\documentclass[11pt,a4paper]{article}
\makeatletter
\newif\ifeleve
\elevefalse

\newif\iflycee@palatino
\lycee@palatinotrue

\RequirePackage[utf8]{inputenc}
\RequirePackage[T1]{fontenc}
\RequirePackage[french]{babel}
\RequirePackage{etoolbox}

\newcommand{\ShorthandsOff}{\@ifundefined{shorthandoff}{}{\shorthandoff{;:!?}}}
\newcommand{\ShorthandsOn}{\@ifundefined{shorthandon}{}{\shorthandon{;:!?}}}
\ShorthandsOff
\AtBeginDocument{\ShorthandsOn}
\AtBeginEnvironment{tikzpicture}{\ShorthandsOff}

\RequirePackage{geometry}
\RequirePackage{fancyhdr}
\RequirePackage{lastpage}
\RequirePackage{multicol}
\RequirePackage{array}
\RequirePackage{enumitem}
\RequirePackage{xcolor}
\RequirePackage{needspace}

\RequirePackage{amsmath,amssymb}
\RequirePackage{mathpazo}
\DeclareMathSizes{10.95}{10.95}{8}{5.5}
\begingroup\catcode`\_=\active \catcode`\^=\active
\gdef_#1{\sb{\scriptscriptstyle #1}}
\gdef^#1{\sp{\scriptscriptstyle\lycee@moinsepais #1}}
\endgroup
\newcommand\lycee@moins{\mathbin{%
  \setbox\z@\hbox{$\m@th\scriptscriptstyle\lycee@moinsorig$}%
  \dimen@=.55\wd\z@ \advance\dimen@-.5pt
  \hbox to.75\wd\z@{\hss$\m@th\scriptscriptstyle\vcenter{\hbox{%
    \tikz\draw[line width=.5pt,line cap=round](0,0)--(\the\dimen@,0);}}$\hss}}}
\begingroup\lccode`\~=`\- \lowercase{\endgroup\let~}\lycee@moins
\newcommand\lycee@moinsepais{\mathcode`\-="8000 }
\AtBeginDocument{\mathchardef\lycee@moinsorig=\mathcode`\-\relax}
\AtBeginDocument{\catcode`\_=12 \mathcode`\_="8000
  \catcode`\^=12 \mathcode`\^="8000 }
\newcommand\lycee@exposants{%
  \fontdimen13\textfont2=.48\fontdimen6\textfont2
  \fontdimen14\textfont2=.48\fontdimen6\textfont2
  \fontdimen15\textfont2=.48\fontdimen6\textfont2
  \fontdimen13\scriptfont2=.48\fontdimen6\scriptfont2
  \fontdimen14\scriptfont2=.48\fontdimen6\scriptfont2
  \fontdimen15\scriptfont2=.48\fontdimen6\scriptfont2}
\every@math@size\expandafter{\the\every@math@size\lycee@exposants}
\newlength{\retraitcalculs}
\setlength{\retraitcalculs}{2em}

\RequirePackage{tikz}
\usetikzlibrary{arrows.meta,calc,babel,shapes.misc}
\RequirePackage[most]{tcolorbox}
\tcbuselibrary{skins,breakable}

\definecolor{coulDef}{HTML}{1F4E79}
\definecolor{coulProp}{HTML}{7030A0}
\definecolor{coulRem}{HTML}{C55A11}
\definecolor{coulEx}{HTML}{2E7D32}
\definecolor{coulExo}{HTML}{B03A2E}
\definecolor{coulPart}{HTML}{1F4E79}
\definecolor{coulMeth}{HTML}{1B6E4F}
\definecolor{coulCourbe}{HTML}{0B5ED7}
\definecolor{coulCourbeB}{HTML}{C00000}
\definecolor{coulCourbeC}{HTML}{2E7D32}
\definecolor{coulGrille}{HTML}{8C8C8C}
\definecolor{coulRep}{HTML}{1B6E4F}
\definecolor{coulBar}{HTML}{2E7D32}

\newcommand{\R}{\mathbb{R}}
\newcommand{\Rs}{\mathbb{R}^{*}}
\renewcommand{\le}{\leqslant}
\renewcommand{\ge}{\geqslant}
\renewcommand{\approx}{\simeq}
\newcommand{\vect}[1]{\overrightarrow{#1}}
\newcommand{\norme}[1]{\left\lVert #1 \right\rVert}
\newcommand{\intervalle}[2]{\left[#1\,;#2\right]}
\RequirePackage{cancel}
\renewcommand{\CancelColor}{\color{coulExo}}
\newcommand{\fois}[1]{\times{\color{coulExo}#1}}
\newcommand{\barre}[1]{\cancel{#1}}

\tikzset{grille/.style={coulGrille,line width=0.4pt}}
\newcommand{\quadrillage}[4]{\draw[grille,step=1] (#1,#2) grid (#3,#4);}
\tikzset{
  croix/.style={cross out,draw,line width=0.6pt,inner sep=0pt,minimum size=4pt},
  croixdroite/.style={croix,rotate=45,line width=0.8pt},
}
\newcommand{\pt}[3]{\path (#1) node[croix]{} node[#2,font=\small,inner sep=2.5pt]{$#3$};}

\newcommand{\lycee@repere}[4]{%
  \draw[grille,step=1] (#1,#2) grid (#3,#4);
  \draw[-{Stealth[length=2mm]},thick] (#1,0)--({#3+0.4},0) node[below left=-1pt]{$x$};
  \draw[-{Stealth[length=2mm]},thick] (0,#2)--(0,{#4+0.4}) node[below left=-1pt]{$y$};
  \node[below left,font=\scriptsize,inner sep=1.5pt] at (0,0) {$0$};}
\newcommand{\lycee@gradx}[1]{\foreach \k/\l in {#1}{%
  \draw (\k,0.07)--(\k,-0.07) node[below,font=\scriptsize,inner sep=1.5pt]{$\l$};}}
\newcommand{\lycee@grady}[1]{\foreach \k/\l in {#1}{%
  \draw (0.07,\k)--(-0.07,\k) node[left,font=\scriptsize,inner sep=1.5pt]{$\l$};}}
\newcommand{\lycee@intff}[2]{\left[#1\,;#2\right]}
\newcommand{\lycee@intfo}[2]{\left[#1\,;#2\right[}
\newcommand{\lycee@intof}[2]{\left]#1\,;#2\right]}
\newcommand{\lycee@intoo}[2]{\left]#1\,;#2\right[}
\newcommand{\lycee@ens}[1]{\left\{#1\right\}}
\AtBeginDocument{%
  \providecommand{\repere}{\lycee@repere}%
  \providecommand{\gradx}{\lycee@gradx}%
  \providecommand{\grady}{\lycee@grady}%
  \providecommand{\Cf}{\mathcal{C}\sb{\scriptscriptstyle f}}%
  \providecommand{\Cg}{\mathcal{C}\sb{\scriptscriptstyle g}}%
  \providecommand{\intff}{\lycee@intff}%
  \providecommand{\intfo}{\lycee@intfo}%
  \providecommand{\intof}{\lycee@intof}%
  \providecommand{\intoo}{\lycee@intoo}%
  \providecommand{\ens}{\lycee@ens}}

\newlist{questions}{enumerate}{3}
\setlist[questions,1]{label=\arabic*\textdegree),leftmargin=*,itemsep=2pt,topsep=3pt}
\setlist[questions,2]{label=\alph*),leftmargin=*,itemsep=1pt,topsep=2pt}
\setlist[questions,3]{label=\roman*),leftmargin=*,itemsep=1pt,topsep=1pt}
\setlist[itemize]{label=\textbullet,leftmargin=*,itemsep=2pt,topsep=3pt}

\newcommand{\besoinplace}[1]{\par\penalty\@M\begingroup
  \dimen@=#1\relax\dimen@ii=\pagegoal\advance\dimen@ii-\pagetotal
  \ifdim\dimen@>\dimen@ii\ifdim\dimen@ii>\z@\vfil\fi\break\fi\endgroup}

\newcommand{\lycee@pied}{\thepage/\pageref{LastPage}}
\newcommand{\entete}[2]{%
  \pagestyle{fancy}\fancyhf{}%
  \fancyhead[L]{\footnotesize\color{coulPart}\textbf{#1}}%
  \fancyhead[R]{\footnotesize\color{coulPart}\textbf{#2}}%
  \fancyfoot[C]{\footnotesize\lycee@pied}%
  \renewcommand{\headrulewidth}{0.4pt}%
  \renewcommand{\headrule}{\hbox to\headwidth{%
    \color{coulPart!40}\leaders\hrule height \headrulewidth\hfill}}}

\RequirePackage{ccicons}
\newcommand{\lycee@licence}{{\scriptsize\color{gray}%
  \copyright\ Luc Destombe --- \ccbyncsa\ CC BY-NC-SA 4.0}}
\newif\iflycee@licence \lycee@licencetrue
\newcommand{\sanslicence}{\lycee@licencefalse}
\AtBeginDocument{\iflycee@licence\AddToHookNext{shipout/foreground}{%
  \put(0,0){\rlap{\kern\dimexpr1in+\hoffset+\oddsidemargin+\textwidth\relax
    \llap{\raisebox{-\dimexpr1in+\voffset+\topmargin+\headheight+\headsep
      +\textheight+\footskip\relax}{\lycee@licence}}}}}\fi}

\newcommand{\formulesserrees}{%
  \setlength{\abovedisplayskip}{3pt}\setlength{\belowdisplayskip}{3pt}%
  \setlength{\abovedisplayshortskip}{2pt}\setlength{\belowdisplayshortskip}{3pt}}

\setlength{\parindent}{0pt}

\RequirePackage[implicit=false,hidelinks,pdfencoding=auto,psdextra,bookmarksopen,
  bookmarksopenlevel=1,pdfcreator={},pdfproducer={}]{hyperref}
\RequirePackage{bookmark}
\hypersetup{pdfauthor={Luc Destombe},pdfkeywords={CC BY-NC-SA 4.0}}
\newcounter{lycee@signet}
\newcommand{\signet}[3][]{\stepcounter{lycee@signet}%
  \edef\lycee@dest{\if\relax\detokenize{#1}\relax
    signet.\arabic{lycee@signet}\else#1\fi}%
  \hypertarget{\lycee@dest}{}\bookmark[level=#2,dest=\lycee@dest]{#3}}
\pdfstringdefDisableCommands{%
  \def\R{R}\def\vect#1{#1}\def\Cf{Cf}\def\Cg{Cg}\def\fois#1{×#1}%
  \def\barre#1{#1}\def\intervalle#1#2{[#1;#2]}\def\intff#1#2{[#1;#2]}%
  \def\textsuperscript#1{#1}\def\dfrac#1#2{#1/#2}\def\frac#1#2{#1/#2}%
  \def\vec#1{#1⃗}}%
\geometry{top=2cm,bottom=1cm,left=1cm,right=1cm,headheight=14pt,footskip=0.6cm}

\RequirePackage{pgfplots}
\pgfplotsset{compat=1.18}
\RequirePackage{tkz-tab}

\newcounter{partiecours}
\renewcommand{\thepartiecours}{\Roman{partiecours}}
\newcommand{\partiecours}[1]{%
  \refstepcounter{partiecours}%
  \Needspace*{8\baselineskip}%
  \bigskip
  \noindent\begin{tcolorbox}[enhanced,colback=coulPart,colframe=coulPart,
      boxrule=0pt,arc=2pt,left=8pt,right=8pt,top=4pt,bottom=4pt,
      before skip=6pt,after skip=10pt]
    \color{white}\bfseries\large \thepartiecours{} --- #1%
    \signet[partie-\arabic{partiecours}]{1}{\thepartiecours{} --- #1}
  \end{tcolorbox}%
  \addcontentsline{toc}{section}{\thepartiecours{} --- #1}%
}

\newtcolorbox{boitecours}[3][]{%
  enhanced,breakable,
  colback=white,colframe=#2,coltitle=white,
  fonttitle=\bfseries,
  attach boxed title to top left={xshift=8pt,yshift=-\tcboxedtitleheight/2},
  boxed title style={colback=#2,arc=2pt,boxrule=0pt},
  boxrule=0.9pt,arc=2pt,
  left=8pt,right=8pt,top=8pt,bottom=6pt,
  title={#3},#1}

\newcounter{methode}
\newenvironment{definitionbox}[1][Définition]%
  {\begin{boitecours}[phantom={\signet{2}{#1}}]{coulDef}{#1}}{\end{boitecours}}
\newenvironment{proprietebox}[1][Propriété]%
  {\begin{boitecours}[phantom={\signet{2}{#1}}]{coulProp}{#1}}{\end{boitecours}}
\newenvironment{methodebox}[1][Méthode]{\stepcounter{methode}%
  \begin{boitecours}[phantom={\signet[methode-\arabic{methode}]{2}{#1}}]{coulMeth}{#1}}%
  {\end{boitecours}}

\newcommand{\etape}[1]{\par\smallskip\textbf{\color{coulMeth}#1}\par\nobreak\smallskip}
\newcommand{\enonce}[1]{\emph{#1}\par\smallskip}
\newcommand{\reponse}[1]{\par\smallskip
  {\footnotesize\itshape\color{black!55}Réponses : #1}}
\newcommand{\demo}[1]{\textbf{\color{coulMeth}Démonstration.} #1}
\newcommand{\afaire}[1]{%
  \par\smallskip
  \noindent{\small\color{coulExo}$\blacktriangleright$\ \textbf{À faire :}
    fiche d'exercices du chapitre, #1.}\par\smallskip}

\newenvironment{calculs}{\setlength{\jot}{5pt}\@fleqntrue\@mathmargin\retraitcalculs
  \setlength{\abovedisplayskip}{4pt}\setlength{\belowdisplayskip}{4pt}%
  \setlength{\abovedisplayshortskip}{4pt}\setlength{\belowdisplayshortskip}{4pt}%
  \csname alignat*\endcsname{2}}%
  {\csname endalignat*\endcsname}
\newcommand{\rem}[1]{\text{\small\itshape\color{black!60}#1}}
\newcommand{\explique}[2]{\par\smallskip\noindent
  \begin{minipage}[t]{0.57\linewidth}\raggedright #1\end{minipage}\hfill
  \begin{minipage}[t]{0.40\linewidth}\raggedright\leavevmode
    {\small\itshape\color{black!60}#2}\end{minipage}\par}
\newcommand{\cas}[1]{\par\smallskip\noindent\textbf{\color{coulExo}#1}\nobreak}
\newcommand{\calc}[1]{\par\smallskip\textbf{\color{coulExo}#1}\par\nobreak}

\newtcolorbox{remarquebox}[1][Remarque]{%
  enhanced,breakable,colback=white,colframe=coulRem,
  boxrule=0pt,leftrule=2.5pt,arc=0pt,outer arc=0pt,
  left=8pt,right=6pt,top=5pt,bottom=5pt,
  before upper={\textbf{\color{coulRem}#1}\par\smallskip}}

\newtcolorbox{exemplebox}[1][Exemple]{%
  enhanced,breakable,colback=white,colframe=coulEx,
  boxrule=0pt,leftrule=2.5pt,arc=0pt,outer arc=0pt,
  left=8pt,right=6pt,top=5pt,bottom=5pt,
  before upper={\textbf{\color{coulEx}#1}\par\smallskip}}

\pgfplotsset{
  repere/.style={
    axis lines=middle,
    axis line style={-{Stealth[length=2.4mm]},thick},
    every axis x label/.style={at={(ticklabel* cs:1.0)},anchor=west},
    every axis y label/.style={at={(ticklabel* cs:1.0)},anchor=south},
    label style={font=\footnotesize},
    tick label style={font=\scriptsize},
    grid=both,
    major grid style={grille},
    minor tick num=0,
    tick align=inside,
    clip mode=individual,
  },
}
\tikzset{
  courbe/.style={coulCourbe,line width=1.1pt,smooth},
  courbeB/.style={coulCourbeB,line width=1.1pt,smooth},
  courbeC/.style={coulCourbeC,line width=1.1pt,smooth},
}

\newlist{etapes}{enumerate}{1}
\setlist[etapes]{label=\textbf{\arabic*.},leftmargin=*,itemsep=1pt,topsep=2pt}

\newlist{expressions}{enumerate}{1}
\setlist[expressions]{label={\Alph*\ $=$},leftmargin=*,itemsep=3pt,topsep=3pt}

\newcounter{exo}
\renewcommand{\theexo}{\ifnum\value{exo}<10 0\fi\arabic{exo}}
\newtcolorbox{exercice}[1][]{%
  enhanced,breakable,
  colback=white,colframe=coulExo,
  boxrule=0.9pt,arc=2pt,
  left=8pt,right=8pt,top=8pt,bottom=6pt,
  coltitle=white,fonttitle=\bfseries,
  attach boxed title to top left={xshift=8pt,yshift=-\tcboxedtitleheight/2},
  boxed title style={colback=coulExo,arc=2pt,boxrule=0pt},
  phantom={\signet{2}{Exercice \arabic{exo}}},
  title={Exercice \theexo\ifx\relax#1\relax\else{} --- #1\fi}}
\newenvironment{exo}[1][\relax]{\refstepcounter{exo}\begin{exercice}[#1]}{\end{exercice}}

  \newenvironment{correction}{\par}{\par}
  \newcommand{\autableau}{}
  \newcommand{\etapeexemples}{\etape{Exemples corrigés}}
  \newcommand{\etapetableau}[1]{\etape{#1}}
  \newenvironment{exempletableau}[1][Exemple]%
    {\begin{exemplebox}[#1]}{\end{exemplebox}}

\RequirePackage{comment}
\excludecomment{correctionavj}

\newcommand{\coursEntete}{}
\newcommand{\coursNiveau}{}
\newcommand{\coursTitre}{}
\newcommand{\coursSurTitre}{}
\newcommand{\coursSousTitre}{}

\newcommand{\enteteCours}[1]{\renewcommand{\coursEntete}{#1}}
\newcommand{\chapitrecours}[2]{\enteteCours{Chapitre #1 --- #2}}
\newcommand{\niveaucours}[1]{\renewcommand{\coursNiveau}{#1}}
\newcommand{\titrecours}[3][]{%
  \renewcommand{\coursTitre}{#2}%
  \renewcommand{\coursSurTitre}{#1}%
  \renewcommand{\coursSousTitre}{#3}}

\pagestyle{fancy}
\fancyhf{}
\fancyhead[L]{\footnotesize\color{coulPart}\textbf{\coursEntete}}
\fancyhead[R]{\footnotesize\color{coulPart}\textbf{\coursNiveau}}
\fancyfoot[C]{\footnotesize\thepage}
\renewcommand{\headrulewidth}{0.4pt}
\renewcommand{\headrule}{\hbox to\headwidth{\color{coulPart!40}\leaders\hrule height \headrulewidth\hfill}}

\setlength{\parskip}{2pt}

\AtBeginDocument{%
  \begin{center}
    \begin{tcolorbox}[enhanced,width=0.72\linewidth,colback=white,colframe=coulPart,
        boxrule=1.6pt,arc=3pt,halign=center,top=10pt,bottom=10pt]
      \ifx\coursSurTitre\empty
        {\Huge\bfseries\color{coulPart} \coursTitre}\\[4pt]
      \else
        {\Huge\bfseries\color{coulPart} \coursTitre}\\[3pt]
        {\Large\bfseries\color{coulPart} \coursSurTitre}\\[5pt]
      \fi
      {\normalsize \coursSousTitre}%
      
    \end{tcolorbox}
  \end{center}%
}
\makeatother

\chapitrecours{7}{Suites arithmétiques et géométriques}
\niveaucours{1\textsuperscript{re} spécialité mathématiques}
\titrecours{SUITES ARITHMÉTIQUES ET GÉOMÉTRIQUES}{Cours --- Première spécialité}

\begin{document}

\newtcblisting{python}{%
  listing only,enhanced,breakable,
  colback=black!3,colframe=black!35,boxrule=0.5pt,arc=1pt,
  left=6pt,right=6pt,top=2pt,bottom=2pt,
  listing options={language=Python,basicstyle=\ttfamily\small,
    keywordstyle=\color{coulMeth}\bfseries,columns=fullflexible,
    keepspaces=true,showstringspaces=false,aboveskip=0pt,belowskip=0pt}}
\newcommand{\reperesuite}[4]{%
  \draw[grille,step=1] (-1,#2) grid (#1,#3);
  \draw[-{Stealth[length=2mm]},thick] (-1,0)--({#1+0.4},0)
    node[below left=-1pt,font=\footnotesize]{$n$};
  \draw[-{Stealth[length=2mm]},thick] (0,#2)--(0,{#3+0.4})
    node[below left=-1pt,font=\footnotesize]{$#4$};
  \node[below left,font=\scriptsize,inner sep=1.5pt] at (0,0) {$0$};}

\begin{remarquebox}[Comment utiliser ce cours]
Chaque partie commence par ce qu'il faut \textbf{savoir} (définitions, propriétés), puis
vient ce qu'il faut \textbf{savoir faire} : une méthode, avec des
\textbf{exemples corrigés} faits en classe, puis des exercices
\og\textbf{À vous de jouer}\fg{} à chercher seul.

\end{remarquebox}

\partiecours{Suites arithmétiques}

\begin{definitionbox}[Définition 1 --- Suite arithmétique]
\begin{minipage}[c]{0.58\linewidth}\raggedright
Une suite $(u_{n})$ est \textbf{arithmétique} si l'on passe de chaque terme au suivant
en \textbf{ajoutant} toujours le même nombre $r$ :
\[
  u_{n+1}=u_{n}+r \quad\text{pour tout entier naturel } n.
\]
Le nombre $r$ est la \textbf{raison} de la suite ; il peut être négatif.
\end{minipage}\hfill
\begin{minipage}[c]{0.38\linewidth}\centering
\begin{tikzpicture}[>={Stealth[length=2mm]}]
  \node (a) at (0,0) {$3$};
  \node (b) at (1.4,0) {$8$};
  \node (c) at (2.8,0) {$13$};
  \node (d) at (4.2,0) {$18$};
  \foreach \x/\y in {a/b,b/c,c/d}
    \draw[->,coulMeth,thick] (\x.north east) to[bend left=45]
      node[above,font=\scriptsize]{$+5$} (\y.north west);
\end{tikzpicture}

\smallskip
{\footnotesize $u_{0}=3$ et raison $r=5$}
\end{minipage}
\end{definitionbox}

\begin{methodebox}[Méthode 1 --- Reconnaître une suite arithmétique et trouver sa raison]
\etapeexemples
\begin{questions}[label=\alph*)]
  \item Une suite commence par $4$ ;\; $7$ ;\; $10$ ;\; $13$ ;\; $16$. Peut-elle être
        arithmétique ? Si oui, donner sa raison.
  \item Même question pour une suite qui commence par $2$ ;\; $5$ ;\; $11$ ;\; $20$.
  \item Une plante mesure $12$~cm le 1\textsuperscript{er} avril et grandit de $1{,}5$~cm
        par semaine. On note $u_{n}$ sa hauteur, en cm, au bout de $n$ semaines. Donner
        $u_{0}$, la relation de récurrence et la raison.
\end{questions}
\begin{correction}
\cas{a)}
\explique{$7-4=3$ ;\; $10-7=3$ ;\; $13-10=3$ ;\; $16-13=3$.}{on calcule l'écart entre
deux termes qui se suivent}
\explique{L'écart est toujours $3$ : la suite peut être arithmétique, de raison
$r=3$.}{on ajoute $3$ à chaque fois}
\cas{b)}
\explique{$5-2=3$ mais $11-5=6$.}{deux écarts différents suffisent}
\explique{La suite n'est pas arithmétique.}{l'écart n'est pas constant}
\cas{c)}
\explique{$u_{0}=12$ et $u_{n+1}=u_{n}+1{,}5$ : la suite est arithmétique, de raison
$r=1{,}5$.}{chaque semaine, on ajoute $1{,}5$~cm}
\end{correction}

\etape{À vous de jouer}
\begin{questions}[label=\alph*)]
  \item Une suite commence par $25$ ;\; $21$ ;\; $17$ ;\; $13$ ;\; $9$. Peut-elle être
        arithmétique ? Si oui, donner sa raison.
  \item Même question pour une suite qui commence par $1$ ;\; $3$ ;\; $9$ ;\; $27$.
  \item Une citerne contient $800$~L d'eau et l'on en utilise $35$~L par jour. On note
        $v_{n}$ le volume, en litres, au bout de $n$ jours. Donner $v_{0}$, la relation
        de récurrence et la raison.
\end{questions}
\reponse{a) Écart toujours $-4$ : arithmétique, de raison $-4$.\; b) Écarts $2$ puis $6$ :
pas arithmétique.\; c) $v_{0}=800$ et $v_{n+1}=v_{n}-35$ : raison $-35$.}
\end{methodebox}

\afaire{exercices 1 à 4}

\partiecours{Terme général d'une suite arithmétique}

\begin{proprietebox}[Propriété 1 --- Terme général]
Si $(u_{n})$ est arithmétique de raison $r$, alors pour tout entier naturel $n$ :
\[
  u_{n}=u_{0}+n\,r
  \qquad\text{et, plus généralement,}\qquad
  u_{n}=u_{p}+(n-p)\,r .
\]
Pour aller de $u_{p}$ à $u_{n}$, on ajoute $(n-p)$ fois la raison.
\end{proprietebox}

\noindent\textbf{Démonstration.}
\begin{calculs}
u_{1} &= u_{0}+r &\qquad& \rem{définition, avec $n=0$}\\
u_{2} &= u_{1}+r = u_{0}+2r &\qquad& \rem{on remplace $u_{1}$}\\
u_{3} &= u_{2}+r = u_{0}+3r &&\\
u_{n} &= u_{0}+n\,r &\qquad& \rem{$n$ sauts de $r$ entre $u_{0}$ et $u_{n}$}
\end{calculs}

\begin{methodebox}[Méthode 2 --- Exprimer $u_{n}$ en fonction de $n$]
\etapeexemples
Pour chaque suite arithmétique, exprimer le terme général en fonction de $n$, puis
calculer le terme demandé.
\begin{questions}[label=\alph*)]
  \item $u_{0}=7$ et $r=-4$ ; calculer $u_{20}$.
  \item $v_{1}=5$ et, pour tout $n\ge 1$, $v_{n+1}=v_{n}+3$ ; calculer $v_{50}$.
  \item $w_{3}=10$ et $r=2{,}5$ ; calculer $w_{11}$.
\end{questions}
\begin{correction}
\cas{a)}
\begin{calculs}
u_{n} &= u_{0}+n\,r &\qquad& \rem{la formule}\\
u_{n} &= 7-4n &\qquad& \rem{$u_{0}=7$ et $r=-4$}\\
u_{20} &= 7-4\times 20 = -73 &\qquad& \rem{on remplace $n$ par $20$}
\end{calculs}
\cas{b)}
\begin{calculs}
v_{n} &= v_{1}+(n-1)\,r &\qquad& \rem{le premier terme est $v_{1}$ : $p=1$}\\
v_{n} &= 5+3(n-1) &\qquad& \rem{$r=3$}\\
v_{n} &= 3n+2 &&\\
v_{50} &= 3\times 50+2 = 152 &&
\end{calculs}
\cas{c)}
\begin{calculs}
w_{n} &= w_{3}+(n-3)\,r &\qquad& \rem{on part de $w_{3}$ : $p=3$}\\
w_{n} &= 10+2{,}5(n-3) &&\\
w_{n} &= 2{,}5n+2{,}5 &\qquad& \rem{$10-7{,}5=2{,}5$}\\
w_{11} &= 2{,}5\times 11+2{,}5 = 30 &&
\end{calculs}
\end{correction}

\etape{À vous de jouer}
Pour chaque suite arithmétique, exprimer le terme général en fonction de $n$, puis
calculer le terme demandé.
\begin{questions}[label=\alph*)]
  \item $a_{0}=-2$ et $r=5$ ; calculer $a_{20}$.
  \item $b_{1}=8$ et, pour tout $n\ge 1$, $b_{n+1}=b_{n}-2$ ; calculer $b_{50}$.
  \item $c_{4}=6$ et $r=1{,}5$ ; calculer $c_{12}$.
\end{questions}
\reponse{a) $a_{n}=5n-2$ et $a_{20}=98$.\; b) $b_{n}=10-2n$ et $b_{50}=-90$.\;
c) $c_{n}=1{,}5n$ et $c_{12}=18$.}
\end{methodebox}

\begin{methodebox}[Méthode 3 --- Trouver la raison à partir de deux termes]
\etapeexemples
Pour chaque suite arithmétique, trouver la raison, le premier terme $u_{0}$ (ou
$v_{0}$), puis le terme général.
\begin{questions}[label=\alph*)]
  \item $u_{4}=11$ et $u_{10}=29$.
  \item $v_{3}=20$ et $v_{8}=5$.
\end{questions}
\begin{correction}
\cas{a)}
\begin{calculs}
u_{10} &= u_{4}+6r &\qquad& \rem{de $u_{4}$ à $u_{10}$ : $6$ sauts de $r$}\\
29 &= 11+6r &&\\
6r &= 18 &&\\
r &= 3 &&
\end{calculs}
\begin{calculs}
u_{0} &= u_{4}-4r &\qquad& \rem{on recule de $4$ sauts}\\
u_{0} &= 11-12 = -1 &&\\
u_{n} &= -1+3n &\qquad& \rem{$u_{n}=u_{0}+n\,r$}
\end{calculs}
\cas{b)}
\begin{calculs}
v_{8} &= v_{3}+5r &\qquad& \rem{$5$ sauts de $r$}\\
5 &= 20+5r &&\\
5r &= -15 &&\\
r &= -3 &&\\[4pt]
v_{0} &= v_{3}-3r &\qquad& \rem{on recule de $3$ sauts}\\
v_{0} &= 20+9 = 29 &&\\
v_{n} &= 29-3n &&
\end{calculs}
\end{correction}

\etape{À vous de jouer}
Pour chaque suite arithmétique, trouver la raison, le premier terme $a_{0}$ (ou
$b_{0}$), puis le terme général.
\begin{questions}[label=\alph*)]
  \item $a_{2}=9$ et $a_{7}=34$.
  \item $b_{5}=12$ et $b_{11}=-6$.
\end{questions}
\reponse{a) $r=5$ ; $a_{0}=-1$ ; $a_{n}=5n-1$.\; b) $r=-3$ ; $b_{0}=27$ ;
$b_{n}=27-3n$.}
\end{methodebox}

\afaire{exercices 5 à 8}

\partiecours{Montrer qu'une suite est (ou n'est pas) arithmétique}

\begin{proprietebox}[Propriété 2 --- Critère]
$(u_{n})$ est arithmétique si et seulement si, pour tout entier naturel $n$,
\[u_{n+1}-u_{n}=r,\]
où $r$ est une \textbf{constante} : un nombre qui ne dépend pas de $n$. Ce nombre $r$
est la raison.

\smallskip
Pour montrer qu'une suite \textbf{n'est pas} arithmétique, il suffit de trouver deux
écarts différents entre des termes qui se suivent.
\end{proprietebox}

\begin{methodebox}[Méthode 4 --- Montrer qu'une suite est (ou n'est pas) arithmétique]
\etapeexemples
Les suites suivantes, définies pour tout entier naturel $n$, sont-elles arithmétiques ?
\begin{questions}[label=\alph*)]
  \item $u_{n}=4n-9$
  \item $v_{n}=n^{2}-n$
  \item $w_{n}=(n+2)^{2}-n^{2}$
\end{questions}
\begin{correction}
\cas{a)}
\begin{calculs}
u_{n+1}-u_{n} &= 4(n+1)-9-(4n-9) &\qquad& \rem{on calcule la différence}\\
u_{n+1}-u_{n} &= 4n+4-9-4n+9 &&\\
u_{n+1}-u_{n} &= 4 &&
\end{calculs}
\explique{La différence ne dépend pas de $n$ : $(u_{n})$ est arithmétique, de raison
$4$ et de premier terme $u_{0}=-9$.}{elle est constante}
\cas{b)}
\begin{calculs}
v_{n+1}-v_{n} &= (n+1)^{2}-(n+1)-(n^{2}-n) &&\\
v_{n+1}-v_{n} &= n^{2}+2n+1-n-1-n^{2}+n &&\\
v_{n+1}-v_{n} &= 2n &&
\end{calculs}
\explique{La différence dépend de $n$ : $(v_{n})$ n'est pas arithmétique.}{$v_{1}-v_{0}=0$
mais $v_{2}-v_{1}=2$}
\cas{c)}
\begin{calculs}
w_{n} &= n^{2}+4n+4-n^{2} &\qquad& \rem{on développe d'abord}\\
w_{n} &= 4n+4 &&
\end{calculs}
\explique{$w_{n}=4n+4$ est de la forme $an+b$ : $(w_{n})$ est arithmétique, de raison
$4$.}{voir la remarque ci-dessous}
\end{correction}

\etape{À vous de jouer}
Les suites suivantes, définies pour tout entier naturel $n$, sont-elles arithmétiques ?
\begin{questions}[label=\alph*)]
  \item $a_{n}=6-2n$
  \item $b_{n}=n^{2}+3n$
  \item $c_{n}=(n+3)^{2}-n^{2}$
\end{questions}
\reponse{a) Oui, de raison $-2$.\; b) Non : $b_{n+1}-b_{n}=2n+4$.\; c) Oui :
$c_{n}=6n+9$, de raison $6$.}
\end{methodebox}

\begin{remarquebox}[Suites de la forme $u_{n}=an+b$]
Une suite définie par $u_{n}=an+b$ est arithmétique, de raison $a$ et de premier terme
$u_{0}=b$ ; réciproquement, toute suite arithmétique s'écrit $u_{n}=r\,n+u_{0}$.
\end{remarquebox}

\afaire{exercices 9 à 12}

\partiecours{Variations : croissance linéaire}

\begin{proprietebox}[Propriété 3 --- Sens de variation et représentation]
\begin{minipage}[c]{0.58\linewidth}\raggedright
Soit $(u_{n})$ arithmétique de raison $r$.
\begin{itemize}
  \item Si $r>0$, $(u_{n})$ est strictement croissante.
  \item Si $r<0$, $(u_{n})$ est strictement décroissante.
  \item Si $r=0$, $(u_{n})$ est constante.
\end{itemize}
Ses points $(n\,;u_{n})$ sont \textbf{alignés} sur la droite d'équation
$y=r\,x+u_{0}$ : on parle de \textbf{croissance linéaire} (la même quantité est ajoutée
à chaque étape).
\end{minipage}\hfill
\begin{minipage}[c]{0.38\linewidth}\centering
\begin{tikzpicture}[x=0.5cm,y=0.5cm]
  \reperesuite{9}{0}{5}{u_{n}}
  \gradx{1,2,3,4,5,6,7,8} \grady{1,2,3,4}
  \draw[coulCourbe,dashed] (0,4)--(8.5,-0.25);
  \foreach \n in {0,1,...,8}
    \path[coulCourbe] (\n,{4-0.5*\n}) node[croix]{};
\end{tikzpicture}

\smallskip
{\footnotesize $u_{0}=4$ et $r=-0{,}5$ : points sur $y=-0{,}5x+4$}
\end{minipage}
\end{proprietebox}

\noindent\textbf{Démonstration.}
\begin{calculs}
u_{n+1}-u_{n} &= u_{n}+r-u_{n} &\qquad& \rem{définition}\\
u_{n+1}-u_{n} &= r &\qquad& \rem{le signe de $r$ donne le sens de variation (ch.~5)}
\end{calculs}
\explique{Les points sont sur la droite $y=r\,x+u_{0}$, car $u_{n}=r\,n+u_{0}$.}{la
fonction affine associée}

\begin{methodebox}[Méthode 5 --- Étudier et représenter une suite arithmétique]
\etapeexemples
La suite $(u_{n})$ est définie par $u_{0}=1$ et, pour tout entier naturel $n$,
$u_{n+1}=u_{n}+0{,}5$.
\begin{questions}[label=\alph*)]
  \item Donner la nature de la suite et son sens de variation.
  \item Représenter les termes $u_{0}$ à $u_{6}$.
  \item Donner une équation de la droite qui porte ces points.
\end{questions}
\begin{correction}
\cas{a)}
\explique{On ajoute $0{,}5$ à chaque terme : $(u_{n})$ est arithmétique de raison
$r=0{,}5$.}{définition}
\explique{$r>0$ : la suite est strictement croissante.}{propriété 3}
\cas{b) et c)}\par\nopagebreak
\begin{minipage}[c]{0.58\linewidth}\raggedright
$u_{n}=1+0{,}5n$ : $u_{0}=1$ ;\; $u_{1}=1{,}5$ ;\; $u_{2}=2$ \dots\; $u_{6}=4$.

\smallskip
Les points sont alignés sur la droite d'équation $y=0{,}5x+1$.
{\small\itshape\color{black!60}coefficient directeur $r$, ordonnée à l'origine $u_{0}$}
\end{minipage}\hfill
\begin{minipage}[c]{0.38\linewidth}\centering
\begin{tikzpicture}[x=0.6cm,y=0.6cm]
  \reperesuite{7}{0}{5}{u_{n}}
  \gradx{1,2,3,4,5,6} \grady{1,2,3,4}
  \draw[coulCourbe,dashed] (0,1)--(6.8,4.4);
  \foreach \n in {0,1,...,6}
    \path[coulCourbe] (\n,{1+0.5*\n}) node[croix]{};
\end{tikzpicture}
\end{minipage}
\end{correction}

\etape{À vous de jouer}
La suite $(v_{n})$ est définie par $v_{0}=6$ et, pour tout entier naturel $n$,
$v_{n+1}=v_{n}-1$.
\begin{questions}[label=\alph*)]
  \item Donner la nature de la suite et son sens de variation.
  \item Représenter les termes $v_{0}$ à $v_{6}$.
  \item Donner une équation de la droite qui porte ces points.
\end{questions}
\reponse{a) Arithmétique de raison $-1$ : strictement décroissante.\; b) Points
$(0\,;6)$, $(1\,;5)$ \dots\; $(6\,;0)$.\; c) $y=-x+6$.}
\end{methodebox}

\begin{methodebox}[Méthode 6 --- Trouver le rang à partir duquel un terme dépasse une valeur]
\etapeexemples
\begin{questions}[label=\alph*)]
  \item Soit $u_{n}=3n-8$. À partir de quel rang a-t-on $u_{n}>100$ ?
  \item Soit $v_{n}=50-4n$. À partir de quel rang a-t-on $v_{n}<0$ ?
\end{questions}
\begin{correction}
\cas{a)}
\begin{calculs}
3n-8 &> 100 &\qquad& \rem{on résout l'inéquation}\\
3n &> 108 &&\\
n &> 36 &&
\end{calculs}
\explique{$n$ est entier : $u_{n}>100$ à partir du rang $37$.}{le premier entier
strictement supérieur à $36$}
\cas{b)}
\begin{calculs}
50-4n &< 0 &&\\
-4n &< -50 &&\\
n &> 12{,}5 &\qquad& \rem{on divise par $-4<0$ : le sens change}
\end{calculs}
\explique{$v_{n}<0$ à partir du rang $13$.}{premier entier après $12{,}5$}
\end{correction}

\etape{À vous de jouer}
\begin{questions}[label=\alph*)]
  \item Soit $a_{n}=5n+7$. À partir de quel rang a-t-on $a_{n}>200$ ?
  \item Soit $b_{n}=80-6n$. À partir de quel rang a-t-on $b_{n}<0$ ?
\end{questions}
\reponse{a) $n>38{,}6$ : à partir du rang $39$.\; b) $n>13{,}3$ environ : à partir du
rang $14$.}
\end{methodebox}

\afaire{exercices 13 à 16}

\partiecours{Somme de termes d'une suite arithmétique}

\noindent On veut calculer vite $S=1+2+3+4+5+6$. On écrit la somme dans un sens, puis
dans l'autre, et on additionne colonne par colonne :
\par\smallskip\noindent\hspace*{\retraitcalculs}$\begin{array}{rcccccccccccc}
S &=& 1 &+& 2 &+& 3 &+& 4 &+& 5 &+& 6\\
S &=& 6 &+& 5 &+& 4 &+& 3 &+& 2 &+& 1\\\hline
2S &=& 7 &+& 7 &+& 7 &+& 7 &+& 7 &+& 7
\end{array}$\par\smallskip
Dans chaque colonne, le nombre du haut augmente de $1$ quand celui du bas diminue de $1$ :
la somme reste égale à $7$. Ainsi $2S=6\times 7=42$, donc $S=21$.

\begin{proprietebox}[Propriété 4 --- Somme des premiers entiers]
Pour tout entier $n\ge 1$ :
\[
  1+2+3+\dots+n=\frac{n(n+1)}{2}.
\]
Plus généralement, pour des termes consécutifs d'une suite arithmétique :
\[
  S=\text{nombre de termes}\times\frac{\text{premier terme}+\text{dernier terme}}{2}.
\]
De $u_{0}$ à $u_{n}$, il y a $n+1$ termes ; de $u_{p}$ à $u_{n}$, il y en a $n-p+1$.
\end{proprietebox}

\noindent\textbf{Démonstration} (même idée que ci-dessus, avec $n$ colonnes).
\par\smallskip\noindent\hspace*{\retraitcalculs}$\begin{array}{rcccccccccc}
S &=& 1 &+& 2 &+& \cdots &+& (n-1) &+& n\\
S &=& n &+& (n-1) &+& \cdots &+& 2 &+& 1\\\hline
2S &=& (n+1) &+& (n+1) &+& \cdots &+& (n+1) &+& (n+1)
\end{array}$\par\smallskip
\begin{calculs}
2S &= n(n+1) &\qquad& \rem{$n$ colonnes égales à $n+1$}\\
S &= \frac{n(n+1)}{2} &&
\end{calculs}

\begin{methodebox}[Méthode 7 --- Calculer $1+2+\dots+n$]
\etapeexemples
Calculer les sommes suivantes.
\begin{questions}[label=\alph*)]
  \item $S=1+2+3+\dots+100$
  \item $T=1+2+3+\dots+40$
  \item $U=2+4+6+\dots+80$
\end{questions}
\begin{correction}
\cas{a)}
\begin{calculs}
S &= \frac{100\times 101}{2} &\qquad& \rem{$n=100$ dans la formule}\\
S &= 5\,050 &\qquad& \rem{le calcul du jeune Gauss}
\end{calculs}
\cas{b)}
\begin{calculs}
T &= \frac{40\times 41}{2} = 820 &\qquad& \rem{$n=40$}
\end{calculs}
\cas{c)}
\begin{calculs}
U &= 2\times(1+2+3+\dots+40) &\qquad& \rem{on met $2$ en facteur}\\
U &= 2\times 820 = 1\,640 &\qquad& \rem{résultat du b)}
\end{calculs}
\end{correction}

\etape{À vous de jouer}
Calculer les sommes suivantes.
\begin{questions}[label=\alph*)]
  \item $A=1+2+3+\dots+200$
  \item $B=1+2+3+\dots+30$
  \item $C=3+6+9+\dots+90$
\end{questions}
\reponse{a) $A=20\,100$.\; b) $B=465$.\; c) $C=3\times 465=1\,395$.}
\end{methodebox}

\begin{methodebox}[Méthode 8 --- Calculer une somme de termes consécutifs]
\etapeexemples
\begin{questions}[label=\alph*)]
  \item Soit $u_{n}=3n+2$. Calculer $S=u_{0}+u_{1}+\dots+u_{20}$.
  \item Calculer $T=15+19+23+\dots+99$.
\end{questions}
\begin{correction}
\cas{a)}
\explique{De $u_{0}$ à $u_{20}$ : $21$ termes ; $u_{0}=2$ et $u_{20}=62$.}{nombre de
termes, premier et dernier}
\begin{calculs}
S &= 21\times\frac{2+62}{2} &\qquad& \rem{nombre de termes $\times$ moyenne des extrêmes}\\
S &= 21\times 32 = 672 &&
\end{calculs}
\cas{b)}
\explique{Suite arithmétique de raison $4$ : $u_{n}=15+4n$.}{on ajoute $4$ à chaque
fois}
\begin{calculs}
15+4n &= 99 &\qquad& \rem{rang du dernier terme}\\
4n &= 84 &&\\
n &= 21 &&
\end{calculs}
\explique{De $u_{0}$ à $u_{21}$ : $22$ termes.}{$21+1$}
\begin{calculs}
T &= 22\times\frac{15+99}{2} &&\\
T &= 22\times 57 = 1\,254 &&
\end{calculs}
\end{correction}

\etape{À vous de jouer}
\begin{questions}[label=\alph*)]
  \item Soit $a_{n}=2n+5$. Calculer $A=a_{0}+a_{1}+\dots+a_{30}$.
  \item Calculer $B=7+12+17+\dots+102$.
\end{questions}
\reponse{a) $31$ termes : $A=31\times 35=1\,085$.\; b) $20$ termes : $B=20\times
54{,}5=1\,090$.}
\end{methodebox}

\afaire{exercices 17 à 20}

\partiecours{Problèmes et suites arithmétiques}

\begin{methodebox}[Méthode 9 --- Résoudre un problème avec une suite arithmétique]
\etapeexemples
Une commune plante $30$ arbres en 2026 et décide d'en planter chaque année $12$ de plus
que l'année précédente. On note $u_{n}$ le nombre d'arbres plantés l'année $2026+n$.
\begin{questions}[label=\alph*)]
  \item Calculer $u_{1}$ et $u_{2}$.
  \item Quelle est la nature de la suite $(u_{n})$ ? Justifier.
  \item Exprimer $u_{n}$ en fonction de $n$.
  \item En quelle année la commune plantera-t-elle plus de $200$ arbres ?
  \item Combien d'arbres aura-t-elle plantés en tout de 2026 à 2035 ?
\end{questions}
\begin{correction}
\cas{a)}
\explique{$u_{1}=30+12=42$ et $u_{2}=42+12=54$.}{$12$ de plus chaque année}
\cas{b)}
\explique{Chaque année, on ajoute $12$ : $u_{n+1}=u_{n}+12$. La suite est arithmétique,
de raison $12$ et de premier terme $u_{0}=30$.}{on cite la relation}
\cas{c)}
\begin{calculs}
u_{n} &= u_{0}+n\,r &\qquad& \rem{terme général}\\
u_{n} &= 30+12n &&
\end{calculs}
\cas{d)}
\begin{calculs}
30+12n &> 200 &&\\
12n &> 170 &&\\
n &> \frac{170}{12} &\qquad& \rem{$\frac{170}{12}\approx 14{,}2$}
\end{calculs}
\explique{Le premier rang est $n=15$ : en $2026+15=2041$.}{$n$ est entier}
\cas{e)}
\explique{De 2026 à 2035 : $u_{0}$ à $u_{9}$, soit $10$ termes ; $u_{9}=30+108=138$.}{on
compte les années}
\begin{calculs}
S &= 10\times\frac{30+138}{2} &\qquad& \rem{nombre de termes $\times$ moyenne des extrêmes}\\
S &= 10\times 84 = 840 &&
\end{calculs}
\explique{La commune aura planté $840$ arbres.}{on conclut dans le contexte}
\end{correction}

\etape{À vous de jouer}
Pour préparer un marathon, une coureuse parcourt $20$~km la première semaine, puis
$3$~km de plus chaque semaine. On note $u_{n}$ la distance, en km, parcourue la semaine
$n$ (la première semaine correspond à $n=0$).
\begin{questions}[label=\alph*)]
  \item Calculer $u_{1}$ et $u_{2}$.
  \item Quelle est la nature de la suite $(u_{n})$ ? Justifier.
  \item Exprimer $u_{n}$ en fonction de $n$.
  \item À partir de quelle semaine parcourt-elle plus de $60$~km ?
  \item Quelle distance aura-t-elle parcourue en tout pendant les $12$ premières
        semaines ?
\end{questions}
\reponse{a) $23$ et $26$.\; b) Arithmétique de raison $3$ et de premier terme $20$.\;
c) $u_{n}=20+3n$.\; d) $n>13{,}3$ environ : semaine $n=14$.\; e) $12\times 36{,}5=438$~km.}
\end{methodebox}

\afaire{exercices 21 à 24}

\partiecours{Suites géométriques}

\begin{definitionbox}[Définition 2 --- Suite géométrique]
\begin{minipage}[c]{0.58\linewidth}\raggedright
Une suite $(u_{n})$ est \textbf{géométrique} si l'on passe de chaque terme au suivant
en \textbf{multipliant} toujours par le même nombre $q$ :
\[
  u_{n+1}=q\times u_{n} \quad\text{pour tout entier naturel } n.
\]
Le nombre $q$ est la \textbf{raison} de la suite.
\end{minipage}\hfill
\begin{minipage}[c]{0.38\linewidth}\centering
\begin{tikzpicture}[>={Stealth[length=2mm]}]
  \node (a) at (0,0) {$5$};
  \node (b) at (1.4,0) {$10$};
  \node (c) at (2.8,0) {$20$};
  \node (d) at (4.2,0) {$40$};
  \foreach \x/\y in {a/b,b/c,c/d}
    \draw[->,coulMeth,thick] (\x.north east) to[bend left=45]
      node[above,font=\scriptsize]{$\times 2$} (\y.north west);
\end{tikzpicture}

\smallskip
{\footnotesize $u_{0}=5$ et raison $q=2$}
\end{minipage}
\end{definitionbox}

\begin{remarquebox}[Pourcentages]
Augmenter de $t\,\%$, c'est multiplier par $1+\frac{t}{100}$ ; diminuer de $t\,\%$,
c'est multiplier par $1-\frac{t}{100}$. Une évolution de $t\,\%$ chaque année donne donc
une suite géométrique : $+4\,\%$ donne $q=1{,}04$, et $-15\,\%$ donne $q=0{,}85$.
\end{remarquebox}

\begin{methodebox}[Méthode 10 --- Reconnaître une suite géométrique et trouver sa raison]
\etapeexemples
\begin{questions}[label=\alph*)]
  \item Une suite commence par $3$ ;\; $6$ ;\; $12$ ;\; $24$ ;\; $48$. Peut-elle être
        géométrique ? Si oui, donner sa raison.
  \item Même question pour une suite qui commence par $2$ ;\; $6$ ;\; $12$ ;\; $20$.
  \item Une ville compte $40\,000$ habitants en 2026 ; sa population augmente de
        $2\,\%$ par an. On note $u_{n}$ la population en $2026+n$. Donner $u_{0}$, la
        relation de récurrence et la raison.
\end{questions}
\begin{correction}
\cas{a)}
\explique{$\frac{6}{3}=2$ ;\; $\frac{12}{6}=2$ ;\; $\frac{24}{12}=2$ ;\;
$\frac{48}{24}=2$.}{on divise chaque terme par le précédent}
\explique{Le quotient est toujours $2$ : la suite peut être géométrique, de raison
$q=2$.}{on multiplie par $2$ à chaque fois}
\cas{b)}
\explique{$\frac{6}{2}=3$ mais $\frac{12}{6}=2$ : la suite n'est pas
géométrique.}{deux quotients différents suffisent}
\cas{c)}
\explique{$u_{0}=40\,000$ et $u_{n+1}=1{,}02\,u_{n}$ : la suite est géométrique, de
raison $q=1{,}02$.}{$+2\,\%$ : on multiplie par $1{,}02$}
\end{correction}

\etape{À vous de jouer}
\begin{questions}[label=\alph*)]
  \item Une suite commence par $80$ ;\; $40$ ;\; $20$ ;\; $10$ ;\; $5$. Peut-elle être
        géométrique ? Si oui, donner sa raison.
  \item Même question pour une suite qui commence par $1$ ;\; $4$ ;\; $9$ ;\; $16$.
  \item Une voiture vaut $25\,000$~€ en 2026 et perd $15\,\%$ de sa valeur chaque année.
        On note $v_{n}$ sa valeur en $2026+n$. Donner $v_{0}$, la relation de
        récurrence et la raison.
\end{questions}
\reponse{a) Quotient toujours $0{,}5$ : géométrique, de raison $0{,}5$.\; b) Quotients
$4$ puis $2{,}25$ : pas géométrique.\; c) $v_{0}=25\,000$ et $v_{n+1}=0{,}85\,v_{n}$ :
raison $0{,}85$.}
\end{methodebox}

\afaire{exercices 25 à 28}

\partiecours{Terme général d'une suite géométrique}

\begin{proprietebox}[Propriété 5 --- Terme général]
Si $(u_{n})$ est géométrique de raison $q$, alors pour tout entier naturel $n$ :
\[
  u_{n}=u_{0}\times q^{n}
  \qquad\text{et, plus généralement,}\qquad
  u_{n}=u_{p}\times q^{\,n-p} .
\]
Pour aller de $u_{p}$ à $u_{n}$, on multiplie $(n-p)$ fois par la raison.
\end{proprietebox}

\noindent\textbf{Démonstration.}
\begin{calculs}
u_{1} &= u_{0}\times q &\qquad& \rem{définition, avec $n=0$}\\
u_{2} &= u_{1}\times q = u_{0}\times q^{2} &\qquad& \rem{on remplace $u_{1}$}\\
u_{3} &= u_{2}\times q = u_{0}\times q^{3} &&\\
u_{n} &= u_{0}\times q^{n} &\qquad& \rem{$n$ multiplications par $q$}
\end{calculs}

\begin{methodebox}[Méthode 11 --- Exprimer $u_{n}$ en fonction de $n$]
\etapeexemples
Pour chaque suite géométrique, exprimer le terme général en fonction de $n$, puis
calculer le terme demandé.
\begin{questions}[label=\alph*)]
  \item $u_{0}=3$ et $q=2$ ; calculer $u_{10}$.
  \item $v_{1}=5$ et, pour tout $n\ge 1$, $v_{n+1}=3v_{n}$ ; calculer $v_{6}$.
  \item $w_{2}=64$ et $q=0{,}5$ ; calculer $w_{8}$.
\end{questions}
\begin{correction}
\cas{a)}
\begin{calculs}
u_{n} &= u_{0}\times q^{n} &\qquad& \rem{la formule}\\
u_{n} &= 3\times 2^{n} &&\\
u_{10} &= 3\times 2^{10} = 3\,072 &\qquad& \rem{$2^{10}=1\,024$}
\end{calculs}
\cas{b)}
\begin{calculs}
v_{n} &= v_{1}\times q^{\,n-1} &\qquad& \rem{le premier terme est $v_{1}$ : $p=1$}\\
v_{n} &= 5\times 3^{\,n-1} &&\\
v_{6} &= 5\times 3^{5} = 1\,215 &\qquad& \rem{$3^{5}=243$}
\end{calculs}
\cas{c)}
\begin{calculs}
w_{n} &= w_{2}\times q^{\,n-2} &\qquad& \rem{on part de $w_{2}$ : $p=2$}\\
w_{n} &= 64\times 0{,}5^{\,n-2} &&\\
w_{8} &= 64\times 0{,}5^{6} = 1 &\qquad& \rem{$0{,}5^{6}=\frac{1}{64}$}
\end{calculs}
\end{correction}

\etape{À vous de jouer}
Pour chaque suite géométrique, exprimer le terme général en fonction de $n$, puis
calculer le terme demandé.
\begin{questions}[label=\alph*)]
  \item $a_{0}=5$ et $q=3$ ; calculer $a_{6}$.
  \item $b_{1}=4$ et, pour tout $n\ge 1$, $b_{n+1}=2b_{n}$ ; calculer $b_{8}$.
  \item $c_{3}=81$ et $q=\frac{1}{3}$ ; calculer $c_{7}$.
\end{questions}
\reponse{a) $a_{n}=5\times 3^{n}$ et $a_{6}=3\,645$.\; b) $b_{n}=4\times 2^{\,n-1}$ et
$b_{8}=512$.\; c) $c_{n}=81\times\left(\frac{1}{3}\right)^{n-3}$ et $c_{7}=1$.}
\end{methodebox}

\begin{methodebox}[Méthode 12 --- Trouver la raison à partir de deux termes]
\etapeexemples
Pour chaque suite géométrique de raison positive, trouver la raison, le premier terme,
puis le terme général.
\begin{questions}[label=\alph*)]
  \item $u_{2}=20$ et $u_{5}=160$.
  \item $v_{1}=400$ et $v_{3}=100$.
\end{questions}
\begin{correction}
\cas{a)}
\begin{calculs}
u_{5} &= u_{2}\times q^{3} &\qquad& \rem{de $u_{2}$ à $u_{5}$ : $3$ multiplications par $q$}\\
160 &= 20\times q^{3} &&\\
q^{3} &= 8 &\qquad& \rem{on divise par $20$}\\
q &= 2 &\qquad& \rem{le nombre dont le cube vaut $8$}
\end{calculs}
\begin{calculs}
u_{0} &= \frac{u_{2}}{q^{2}} &\qquad& \rem{on recule de $2$ rangs : on divise deux fois par $q$}\\
u_{0} &= \frac{20}{4} = 5 &&\\
u_{n} &= 5\times 2^{n} &&
\end{calculs}
\cas{b)}
\begin{calculs}
v_{3} &= v_{1}\times q^{2} &\qquad& \rem{$2$ multiplications par $q$}\\
100 &= 400\times q^{2} &&\\
q^{2} &= 0{,}25 &&\\
q &= 0{,}5 &\qquad& \rem{$q>0$ et $0{,}5^{2}=0{,}25$}\\[4pt]
v_{0} &= \frac{v_{1}}{q} = \frac{400}{0{,}5} &\qquad& \rem{on recule d'un rang}\\
v_{0} &= 800 &&\\
v_{n} &= 800\times 0{,}5^{n} &&
\end{calculs}
\end{correction}

\etape{À vous de jouer}
Pour chaque suite géométrique de raison positive, trouver la raison, le premier terme,
puis le terme général.
\begin{questions}[label=\alph*)]
  \item $a_{2}=18$ et $a_{5}=486$.
  \item $b_{1}=250$ et $b_{3}=10$.
\end{questions}
\reponse{a) $q=3$ ; $a_{0}=2$ ; $a_{n}=2\times 3^{n}$.\; b) $q=0{,}2$ ; $b_{0}=1\,250$ ;
$b_{n}=1\,250\times 0{,}2^{n}$.}
\end{methodebox}

\afaire{exercices 29 à 32}

\partiecours{Montrer qu'une suite est (ou n'est pas) géométrique}

\begin{proprietebox}[Propriété 6 --- Critère]
$(u_{n})$ est géométrique s'il existe un nombre $q$ tel que, pour tout $n$,
$u_{n+1}=q\times u_{n}$ : on exprime $u_{n+1}$ et on fait apparaître $u_{n}$.

\smallskip
Pour montrer qu'une suite \textbf{n'est pas} géométrique, il suffit de trouver deux
quotients différents entre des termes qui se suivent.
\end{proprietebox}

\begin{methodebox}[Méthode 13 --- Montrer qu'une suite est (ou n'est pas) géométrique]
\etapeexemples
Les suites suivantes, définies pour tout entier naturel $n$, sont-elles géométriques ?
\begin{questions}[label=\alph*)]
  \item $u_{n}=7\times 4^{n}$
  \item $v_{n}=2^{n}+1$
  \item $w_{n}=\dfrac{6}{2^{n}}$
\end{questions}
\begin{correction}
\cas{a)}
\begin{calculs}
u_{n+1} &= 7\times 4^{n+1} &\qquad& \rem{on exprime le terme suivant}\\
u_{n+1} &= 4\times 7\times 4^{n} &\qquad& \rem{$4^{n+1}=4\times 4^{n}$}\\
u_{n+1} &= 4\,u_{n} &&
\end{calculs}
\explique{$(u_{n})$ est géométrique, de raison $4$ et de premier terme
$u_{0}=7$.}{$u_{n+1}=q\times u_{n}$}
\cas{b)}
\explique{$v_{0}=2$ ;\; $v_{1}=3$ ;\; $v_{2}=5$ : $\frac{v_{1}}{v_{0}}=1{,}5$ mais
$\frac{v_{2}}{v_{1}}=\frac{5}{3}$.}{un contre-exemple sur les premiers termes}
\explique{Les quotients diffèrent : $(v_{n})$ n'est pas géométrique.}{il suffit d'un
écart}
\cas{c)}
\begin{calculs}
w_{n+1} &= \frac{6}{2^{n+1}} &\qquad& \rem{on exprime le terme suivant}\\
w_{n+1} &= \frac{6}{2\times 2^{n}} &\qquad& \rem{$2^{n+1}=2\times 2^{n}$}\\
w_{n+1} &= 0{,}5\times\frac{6}{2^{n}} = 0{,}5\,w_{n} &&
\end{calculs}
\explique{$(w_{n})$ est géométrique, de raison $0{,}5$ et de premier terme
$w_{0}=6$.}{on reconnaît $w_{n}$}
\end{correction}

\etape{À vous de jouer}
Les suites suivantes, définies pour tout entier naturel $n$, sont-elles géométriques ?
\begin{questions}[label=\alph*)]
  \item $a_{n}=2\times 5^{n}$
  \item $b_{n}=3^{n}+2$
  \item $c_{n}=\dfrac{10}{3^{n}}$
\end{questions}
\reponse{a) Oui, de raison $5$.\; b) Non : $\frac{b_{1}}{b_{0}}=\frac{5}{3}$ mais
$\frac{b_{2}}{b_{1}}=\frac{11}{5}$.\; c) Oui, de raison $\frac{1}{3}$.}
\end{methodebox}

\afaire{exercices 33 à 36}

\partiecours{Variations : croissance exponentielle}

\begin{proprietebox}[Propriété 7 --- Sens de variation]
Soit $(u_{n})$ géométrique de raison $q>0$ et de premier terme $u_{0}>0$.
\begin{itemize}
  \item Si $q>1$, $(u_{n})$ est strictement croissante.
  \item Si $0<q<1$, $(u_{n})$ est strictement décroissante.
  \item Si $q=1$, $(u_{n})$ est constante.
\end{itemize}
Si $u_{0}<0$, le sens est inversé.
\end{proprietebox}

\noindent\textbf{Démonstration.}
\begin{calculs}
u_{n+1}-u_{n} &= u_{0}\,q^{n+1}-u_{0}\,q^{n} &\qquad& \rem{terme général}\\
u_{n+1}-u_{n} &= u_{0}\,q^{n}(q-1) &\qquad& \rem{facteur commun $u_{0}\,q^{n}$}
\end{calculs}
\explique{$u_{0}>0$ et $q^{n}>0$ : la différence a le signe de $q-1$.}{positive si
$q>1$, négative si $q<1$}

\begin{remarquebox}[Croissance linéaire, croissance exponentielle]
\begin{minipage}[c]{0.58\linewidth}\raggedright
Une suite arithmétique ($r>0$) ajoute la \textbf{même quantité} à chaque rang : ses
points sont alignés (croissance \textbf{linéaire}).

\smallskip
Une suite géométrique ($q>1$) augmente du \textbf{même pourcentage} : ses points
montent de plus en plus vite (croissance \textbf{exponentielle}). Tôt ou tard, elle
dépasse n'importe quelle suite arithmétique.
\end{minipage}\hfill
\begin{minipage}[c]{0.38\linewidth}\centering
\begin{tikzpicture}[x=0.5cm,y=0.5cm]
  \reperesuite{5}{0}{9}{y}
  \gradx{1,2,3,4} \grady{1,2,3,4,5,6,7,8}
  \foreach \n in {0,1,...,4}{
    \path[coulCourbe] (\n,{0.5+1.5*\n}) node[croix]{};
    \path[coulCourbeB] (\n,{0.5*2^\n}) node[croixdroite]{};}
  \node[coulCourbe,font=\footnotesize,anchor=west] at (4.3,6.5) {$a_{n}$};
  \node[coulCourbeB,font=\footnotesize,anchor=west] at (4.3,8.2) {$g_{n}$};
\end{tikzpicture}

\smallskip
{\footnotesize $a_{n}=0{,}5+1{,}5n$ (\texttimes) et $g_{n}=0{,}5\times 2^{n}$ (+)}
\end{minipage}
\end{remarquebox}

\begin{methodebox}[Méthode 14 --- Étudier le sens de variation d'une suite géométrique]
\etapeexemples
Étudier le sens de variation de chaque suite, définie pour tout entier naturel $n$.
\begin{questions}[label=\alph*)]
  \item $u_{n}=3\times 1{,}2^{n}$
  \item $v_{0}=500$ et $v_{n+1}=0{,}9\,v_{n}$
  \item $w_{n}=-2\times 3^{n}$
\end{questions}
\begin{correction}
\cas{a)}
\explique{Géométrique de raison $q=1{,}2$ et de premier terme $u_{0}=3$.}{forme
$u_{0}\times q^{n}$}
\explique{$u_{0}>0$ et $q>1$ : $(u_{n})$ est strictement croissante.}{propriété 7}
\cas{b)}
\explique{Géométrique de raison $q=0{,}9$ et de premier terme $v_{0}=500$.}{forme
$v_{n+1}=q\times v_{n}$}
\explique{$v_{0}>0$ et $0<q<1$ : $(v_{n})$ est strictement décroissante.}{propriété 7}
\cas{c)}
\explique{Géométrique de raison $q=3$ et de premier terme $w_{0}=-2$.}{forme
$w_{0}\times q^{n}$}
\explique{$q>1$ mais $w_{0}<0$ : le sens est inversé, $(w_{n})$ est strictement
décroissante.}{$-2$ ; $-6$ ; $-18$\dots}
\end{correction}

\etape{À vous de jouer}
Étudier le sens de variation de chaque suite, définie pour tout entier naturel $n$.
\begin{questions}[label=\alph*)]
  \item $a_{n}=7\times 1{,}5^{n}$
  \item $b_{0}=2\,000$ et $b_{n+1}=0{,}75\,b_{n}$
  \item $c_{n}=-5\times 2^{n}$
\end{questions}
\reponse{a) $a_{0}>0$ et $q=1{,}5>1$ : strictement croissante.\; b) $b_{0}>0$ et
$q=0{,}75<1$ : strictement décroissante.\; c) $c_{0}<0$ et $q=2>1$ : strictement
décroissante.}
\end{methodebox}

\begin{methodebox}[Méthode 15 --- Comparer une croissance linéaire et une croissance exponentielle]
\etapeexemples
Une entreprise propose deux contrats. Contrat A : $1\,800$~€ par mois la première
année, puis $50$~€ de plus chaque année. Contrat B : $1\,700$~€ par mois la première
année, puis une hausse de $4\,\%$ chaque année. On note $a_{n}$ et $b_{n}$ les salaires
mensuels l'année $n$ (la première année correspond à $n=0$).
\begin{questions}[label=\alph*)]
  \item Donner la nature de chaque suite et exprimer $a_{n}$ et $b_{n}$ en fonction de
        $n$.
  \item À l'aide du tableau de la calculatrice, déterminer à partir de quelle année le
        contrat B est le plus avantageux.
\end{questions}
\begin{correction}
\cas{a)}
\explique{$a_{n+1}=a_{n}+50$ : arithmétique de raison $50$, donc
$a_{n}=1\,800+50n$.}{on ajoute $50$~€}
\explique{$b_{n+1}=1{,}04\,b_{n}$ : géométrique de raison $1{,}04$, donc
$b_{n}=1\,700\times 1{,}04^{n}$.}{$+4\,\%$ : on multiplie par $1{,}04$}
\cas{b)}
\explique{On entre les deux suites et on lit le tableau :}{application \texttt{Suites}
(ch.~5)}
\begin{center}
\renewcommand{\arraystretch}{1.3}
\begin{tabular}{|c|c|c|}
\hline
$n$ & $a_{n}$ & $b_{n}$ (arrondi)\\\hline
$4$ & $2\,000$ & $1\,988{,}76$\\\hline
$5$ & $2\,050$ & $2\,068{,}31$\\\hline
\end{tabular}
\end{center}
\explique{$b_{4}<a_{4}$ mais $b_{5}>a_{5}$ : le contrat B devient le plus avantageux à
partir de l'année $n=5$, la sixième année.}{puis l'écart ne cesse de grandir :
croissance exponentielle}
\end{correction}

\etape{À vous de jouer}
Une chaîne vidéo A compte $5\,000$ abonnés et en gagne $300$ par mois. Une chaîne B
compte $3\,000$ abonnés et leur nombre augmente de $6\,\%$ par mois. On note $a_{n}$ et
$b_{n}$ les nombres d'abonnés au bout de $n$ mois.
\begin{questions}[label=\alph*)]
  \item Donner la nature de chaque suite et exprimer $a_{n}$ et $b_{n}$ en fonction de
        $n$.
  \item À l'aide du tableau de la calculatrice, déterminer au bout de combien de mois la
        chaîne B dépasse la chaîne A.
\end{questions}
\reponse{a) $a_{n}=5\,000+300n$ (arithmétique) et $b_{n}=3\,000\times 1{,}06^{n}$
(géométrique).\; b) $b_{24}\approx 12\,147<12\,200=a_{24}$ et
$b_{25}\approx 12\,876>12\,500=a_{25}$ : au bout de $25$ mois.}
\end{methodebox}

\afaire{exercices 37 à 40}

\partiecours{Somme de termes d'une suite géométrique}

\noindent On veut calculer $S=1+2+4+8+16$. On multiplie la somme par la raison $2$ et on
écrit $2S$ sous $S$, décalé d'une colonne : les termes égaux sont en regard. On soustrait
les deux lignes : les termes en regard s'éliminent.
\par\smallskip\noindent\hspace*{\retraitcalculs}$\begin{array}{rcccccccccccc}
S &=& 1 &+& \barre{2} &+& \barre{4} &+& \barre{8} &+& \barre{16} && \\
2S &=& && \barre{2} &+& \barre{4} &+& \barre{8} &+& \barre{16} &+& 32\\\hline
S-2S &=& 1 &&&&&&&&& - & 32
\end{array}$\par\smallskip
\begin{calculs}
-S &= -31 &\qquad& \rem{$S-2S=-S$ et $1-32=-31$}\\
S &= 31 &&
\end{calculs}
Seuls le premier terme et le dernier terme de $2S$ restent : on parle de
\textbf{télescopage}.

\begin{proprietebox}[Propriété 8 --- Somme des puissances de $q$]
Pour tout réel $q\neq 1$ et tout entier naturel $n$ :
\[
  1+q+q^{2}+\dots+q^{n}=\frac{1-q^{n+1}}{1-q}.
\]
Plus généralement, pour des termes consécutifs d'une suite géométrique de raison
$q\neq 1$ :
\[
  S=\text{premier terme}\times\frac{1-q^{\,\text{nombre de termes}}}{1-q}.
\]
\end{proprietebox}

\noindent\textbf{Démonstration} (même idée : on écrit $q\,S$ sous $S$, décalé d'une
colonne, et on soustrait).
\par\smallskip\noindent\hspace*{\retraitcalculs}$\begin{array}{rcccccccccccc}
S &=& 1 &+& \barre{q} &+& \barre{q^{2}} &+& \cdots &+& \barre{q^{n}} && \\
q\,S &=& && \barre{q} &+& \barre{q^{2}} &+& \cdots &+& \barre{q^{n}} &+& q^{n+1}\\\hline
S-q\,S &=& 1 &&&&&&&&& - & q^{n+1}
\end{array}$\par\smallskip
\begin{calculs}
S(1-q) &= 1-q^{n+1} &\qquad& \rem{on factorise par $S$}\\
S &= \frac{1-q^{n+1}}{1-q} &\qquad& \rem{$q\neq 1$ : on divise par $1-q$}
\end{calculs}

\begin{methodebox}[Méthode 16 --- Calculer une somme de termes d'une suite géométrique]
\etapeexemples
Calculer les sommes suivantes.
\begin{questions}[label=\alph*)]
  \item $S=1+3+3^{2}+\dots+3^{6}$
  \item $T=5+10+20+\dots+5\times 2^{9}$
  \item $U=1+0{,}5+0{,}5^{2}+\dots+0{,}5^{9}$ (valeur arrondie au millième)
\end{questions}
\begin{correction}
\cas{a)}
\begin{calculs}
S &= \frac{1-3^{7}}{1-3} &\qquad& \rem{$q=3$ et $n=6$ : exposant $n+1=7$}\\
S &= \frac{1-2\,187}{-2} = 1\,093 &&
\end{calculs}
\cas{b)}
\begin{calculs}
T &= 5\times(1+2+\dots+2^{9}) &\qquad& \rem{premier terme $5$ en facteur}\\
T &= 5\times\frac{1-2^{10}}{1-2} &\qquad& \rem{$10$ termes, de $2^{0}$ à $2^{9}$}\\
T &= 5\times 1\,023 = 5\,115 &&
\end{calculs}
\cas{c)}
\begin{calculs}
U &= \frac{1-0{,}5^{10}}{1-0{,}5} &\qquad& \rem{$q=0{,}5$ et $10$ termes}\\
U &\approx 1{,}998 &\qquad& \rem{à la calculatrice}
\end{calculs}
\end{correction}

\etape{À vous de jouer}
Calculer les sommes suivantes.
\begin{questions}[label=\alph*)]
  \item $A=1+2+2^{2}+\dots+2^{9}$
  \item $B=3+12+48+\dots+3\times 4^{5}$
  \item $C=1+0{,}1+0{,}1^{2}+\dots+0{,}1^{5}$ (valeur arrondie au millième)
\end{questions}
\reponse{a) $A=1\,023$.\; b) $B=4\,095$.\; c) $C\approx 1{,}111$.}
\end{methodebox}

\afaire{exercices 41 à 44}

\partiecours{Problèmes et suites géométriques}

\begin{methodebox}[Méthode 17 --- Résoudre un problème avec une suite géométrique]
\etapeexemples
Une jeune entreprise a vendu $1\,200$ vélos électriques en 2026. Elle prévoit une hausse
de ses ventes de $15\,\%$ par an. On note $u_{n}$ le nombre de vélos vendus l'année
$2026+n$.
\begin{questions}[label=\alph*)]
  \item Calculer $u_{1}$ et $u_{2}$.
  \item Quelle est la nature de la suite $(u_{n})$ ? Préciser sa raison.
  \item Exprimer $u_{n}$ en fonction de $n$.
  \item Combien de vélos l'entreprise vendra-t-elle en 2032 ?
  \item Compléter la fonction ci-dessous pour qu'elle renvoie le nombre d'années au bout
        duquel les ventes dépassent $5\,000$ vélos, puis donner cette année.
  \item Combien de vélos aura-t-elle vendus en tout de 2026 à 2035 ?
\end{questions}
\begin{python}
def seuil():
    u = 1200
    n = 0
    while .......... :
        u = ..........
        n = n + 1
    return n
\end{python}
\begin{correction}
\cas{a)}
\begin{calculs}
u_{1} &= 1{,}15\times 1\,200 = 1\,380 &\qquad& \rem{$+15\,\%$ : on multiplie par $1{,}15$}\\
u_{2} &= 1{,}15\times 1\,380 = 1\,587 &&
\end{calculs}
\cas{b)}
\explique{$u_{n+1}=1{,}15\,u_{n}$ : la suite est géométrique, de raison $1{,}15$ et de
premier terme $u_{0}=1\,200$.}{on cite la relation}
\cas{c)}
\begin{calculs}
u_{n} &= u_{0}\times q^{n} &\qquad& \rem{terme général}\\
u_{n} &= 1\,200\times 1{,}15^{n} &&
\end{calculs}
\cas{d)}
\begin{calculs}
u_{6} &= 1\,200\times 1{,}15^{6} &\qquad& \rem{$2032=2026+6$}\\
u_{6} &\approx 2\,776 &\qquad& \rem{arrondi à l'unité}
\end{calculs}
\cas{e)}\par\nopagebreak
\begin{minipage}[t]{0.42\linewidth}\vspace{0pt}
\begin{python}
def seuil():
    u = 1200
    n = 0
    while u <= 5000 :
        u = 1.15*u
        n = n + 1
    return n
\end{python}
\end{minipage}\hfill
\begin{minipage}[t]{0.54\linewidth}\vspace{0pt}\raggedright
{\small\itshape\color{black!60}tant que les ventes ne dépassent pas $5\,000$, on passe à
l'année suivante (on multiplie par $1{,}15$)}

\smallskip
$u_{10}\approx 4\,855$ et $u_{11}\approx 5\,583$ : \texttt{seuil()} renvoie $11$, les
ventes dépassent $5\,000$ vélos en $2026+11=2037$.
\end{minipage}
\cas{f)}
\explique{De 2026 à 2035 : $u_{0}$ à $u_{9}$, soit $10$ termes.}{on compte les années}
\begin{calculs}
S &= 1\,200\times\frac{1-1{,}15^{10}}{1-1{,}15} &\qquad& \rem{premier terme $\times\frac{1-q^{10}}{1-q}$}\\
S &\approx 24\,364 &\qquad& \rem{à la calculatrice}
\end{calculs}
\explique{L'entreprise aura vendu environ $24\,364$ vélos.}{on conclut dans le contexte}
\end{correction}

\etape{À vous de jouer}
Une application a été téléchargée $5\,000$ fois le premier mois. Ensuite, le nombre de
téléchargements baisse de $10\,\%$ chaque mois. On note $u_{n}$ le nombre de
téléchargements le mois $n$ (le premier mois correspond à $n=0$).
\begin{questions}[label=\alph*)]
  \item Calculer $u_{1}$ et $u_{2}$.
  \item Quelle est la nature de la suite $(u_{n})$ ? Préciser sa raison.
  \item Exprimer $u_{n}$ en fonction de $n$.
  \item Combien de téléchargements y aura-t-il le mois $n=6$ ?
  \item Compléter la fonction ci-dessous pour qu'elle renvoie le premier mois où les
        téléchargements passent sous $1\,000$, puis donner ce mois.
  \item Combien de téléchargements y aura-t-il eu en tout pendant les $12$ premiers
        mois ?
\end{questions}
\begin{python}
def seuil():
    u = 5000
    n = 0
    while .......... :
        u = ..........
        n = n + 1
    return n
\end{python}
\reponse{a) $4\,500$ et $4\,050$.\; b) Géométrique de raison $0{,}9$.\;
c) $u_{n}=5\,000\times 0{,}9^{n}$.\; d) Environ $2\,657$.\; e) \texttt{while u >= 1000}
et \texttt{u = 0.9*u} ; le mois $n=16$.\; f) Environ $35\,879$.}
\end{methodebox}

\afaire{exercices 45 à 48}

\partiecours{Suites arithmético-géométriques}

\begin{definitionbox}[Définition 3 --- Suite arithmético-géométrique]
Une suite définie par $u_{n+1}=a\,u_{n}+b$ (avec $a\neq 1$ et $b\neq 0$) est dite
\textbf{arithmético-géométrique} : on multiplie, puis on ajoute. Elle n'est en général
ni arithmétique ni géométrique (exemple : les abonnés d'une médiathèque, ch.~5).

\smallskip
Pour l'étudier, on utilise une \textbf{suite auxiliaire} $v_{n}=u_{n}-\ell$ (le nombre
$\ell$ est donné par l'énoncé), qui est géométrique : on en déduit $v_{n}$, puis $u_{n}$.
\end{definitionbox}

\begin{methodebox}[Méthode 18 --- Étudier une suite arithmético-géométrique]
\etapeexemples
La suite $(u_{n})$ est définie par $u_{0}=200$ et, pour tout entier naturel $n$,
$u_{n+1}=0{,}5\,u_{n}+60$.
\begin{questions}[label=\alph*)]
  \item Calculer $u_{1}$ et $u_{2}$, puis montrer que $(u_{n})$ n'est ni arithmétique
        ni géométrique.
  \item On pose $v_{n}=u_{n}-120$. Montrer que $(v_{n})$ est géométrique ; préciser sa
        raison et son premier terme.
  \item Exprimer $v_{n}$, puis $u_{n}$, en fonction de $n$.
  \item Étudier le sens de variation de $(u_{n})$.
  \item Conjecturer la limite de $(u_{n})$.
\end{questions}
\begin{correction}
\cas{a)}
\begin{calculs}
u_{1} &= 0{,}5\times 200+60 = 160 &\qquad& \rem{la relation avec $n=0$}\\
u_{2} &= 0{,}5\times 160+60 = 140 &&
\end{calculs}
\explique{$u_{1}-u_{0}=-40$ mais $u_{2}-u_{1}=-20$ : pas arithmétique.}{deux écarts
différents}
\explique{$\frac{u_{1}}{u_{0}}=0{,}8$ mais $\frac{u_{2}}{u_{1}}=0{,}875$ : pas
géométrique.}{deux quotients différents}
\cas{b)}
\begin{calculs}
v_{n+1} &= u_{n+1}-120 &\qquad& \rem{définition de $v_{n}$, au rang $n+1$}\\
v_{n+1} &= 0{,}5\,u_{n}+60-120 &\qquad& \rem{relation de récurrence de $(u_{n})$}\\
v_{n+1} &= 0{,}5\,u_{n}-60 &&\\
v_{n+1} &= 0{,}5(u_{n}-120) &\qquad& \rem{on factorise par $0{,}5$ : $0{,}5\times 120=60$}\\
v_{n+1} &= 0{,}5\,v_{n} &\qquad& \rem{on reconnaît $v_{n}=u_{n}-120$}
\end{calculs}
\begin{calculs}
v_{0} &= u_{0}-120 = 80 &\qquad& \rem{premier terme}
\end{calculs}
\explique{$(v_{n})$ est géométrique, de raison $0{,}5$ et de premier terme
$v_{0}=80$.}{forme $v_{n+1}=q\times v_{n}$}
\cas{c)}
\begin{calculs}
v_{n} &= 80\times 0{,}5^{n} &\qquad& \rem{$v_{0}\times q^{n}$}\\
u_{n} &= v_{n}+120 &\qquad& \rem{car $v_{n}=u_{n}-120$}\\
u_{n} &= 120+80\times 0{,}5^{n} &&
\end{calculs}
\explique{Contrôle : $u_{2}=120+80\times 0{,}25=140$.}{on retrouve le a)}
\cas{d)}
\explique{$(v_{n})$ est géométrique avec $v_{0}>0$ et $0<q<1$ : elle est strictement
décroissante.}{propriété 7}
\explique{$u_{n}=v_{n}+120$ : $(u_{n})$ est aussi strictement décroissante.}{ajouter
$120$ ne change pas le sens}
\cas{e)}
\explique{$0{,}5^{n}$ devient aussi proche de $0$ que l'on veut, donc $u_{n}$ se
rapproche de $120$ : on conjecture $\displaystyle\lim_{n\to+\infty}u_{n}=120$.}{par
exemple $0{,}5^{20}<0{,}000\,001$}
\end{correction}

\etape{À vous de jouer}
La suite $(a_{n})$ est définie par $a_{0}=50$ et, pour tout entier naturel $n$,
$a_{n+1}=0{,}6\,a_{n}+8$.
\begin{questions}[label=\alph*)]
  \item Calculer $a_{1}$ et $a_{2}$, puis montrer que $(a_{n})$ n'est ni arithmétique
        ni géométrique.
  \item On pose $b_{n}=a_{n}-20$. Montrer que $(b_{n})$ est géométrique ; préciser sa
        raison et son premier terme.
  \item Exprimer $b_{n}$, puis $a_{n}$, en fonction de $n$.
  \item Étudier le sens de variation de $(a_{n})$.
  \item Conjecturer la limite de $(a_{n})$.
\end{questions}
\reponse{a) $a_{1}=38$ et $a_{2}=30{,}8$ ; écarts $-12$ et $-7{,}2$ ; quotients $0{,}76$ et
environ $0{,}81$.\; b) $b_{n+1}=0{,}6\,b_{n}$ et $b_{0}=30$.\; c) $b_{n}=30\times
0{,}6^{n}$ et $a_{n}=20+30\times 0{,}6^{n}$.\; d) Strictement décroissante.\; e) $20$.}
\end{methodebox}
\afaire{exercices 49 à 55}

\end{document}
